\documentclass[10pt,a4paper]{amsart}
\usepackage[margin=19mm]{geometry}
\usepackage{amsmath,amssymb,mathtools}
\usepackage[scr=rsfs,cal=euler]{mathalfa}
\usepackage{xfrac}
\usepackage[unicode,hidelinks,bookmarks=false]{hyperref}
\newcommand{\ie}{i.e.\ }
\newcommand{\cf}{cf.\ }
\newcommand{\resp}{respectively\ }
\newcommand{\Subsection}{\subsection}
\providecommand{\nobreakdash}{\nobreak\hbox{-}}
\setlength{\emergencystretch}{2em}
\multlinegap0pt
\newtheorem{proposition}{Proposition}
    \def\c{{\mathbb C}}
    \def\r{{\mathbb R}}
\title{Approximating functions on $\r^+$ by exponential sums}
\author{Alexey Kuznetsov, Armin Mohammadioroojeh}
\date{Source: arXiv:2508.19095v2 (CC BY 4.0)}
\begin{document}
\maketitle

\noindent\emph{Source and licence.} Approximating functions on $\r^+$ by exponential sums. Version arXiv:2508.19095v2 (CC BY 4.0), distributed by arXiv under the Creative Commons Attribution 4.0 International licence, http://creativecommons.org/licenses/by/4.0/, as recorded in the arXiv licence metadata for this exact version and shown on the abstract page https://arxiv.org/abs/2508.19095v2. Full licence notice: \texttt{licenses/arxiv-2508.19095-CC-BY-4.0.txt}.\par\smallskip

\noindent\textit{Editorial scope.} Counted unit: the article's proposition prop\_unique (source0.tex lines 181--183). The article's complete original proof of it is delivered with it (source0.tex lines 185--206, one \texttt{proof} environment of 22 lines). No mathematics is written, repaired or re-derived: the statement and the proof are the article's own lines, in the article's own notation, and no retained line is substituted or re-flowed.  Retained uncounted as the source-local context the proof needs: lines 121--123; lines 158--160; lines 168--170.  Nothing was omitted from any retained span, so the package carries no omission record.  No retained line contains a citation, so the wrapper carries no bibliography and no bibliography entry.  No retained line contains a figure environment, an image-inclusion command or drawing code, so no image is delivered and the package declares no asset dependency.  Editorial formatting only, in addition: an AMS \texttt{amsart} preamble that restates the article's own declarations and macros used by the retained lines, namely the environments \texttt{proposition} on separate counters, together with the standard packages those lines need.  The retained environments print in this document as Proposition 1 in order of appearance, whereas the article prints the same items with the numbering of its own full document.  The complete native source of the article as distributed in the arXiv e-print for this exact version is bundled unchanged as \texttt{sources/183926/source0.tex}.
\begin{equation}\label{eqn:Phi_partial_fraction}
R(z)=\sum_{j=1}^M \frac{c_j}{z+\lambda_j}.
\end{equation}

\begin{equation}\label{eqn:Phi_at_z_j}
R(z_j)=F(z_j), \qquad j=1,2,\dots,p, 
\end{equation}

\begin{equation}\label{eqn:Phi_at_infty}
R(z)=\sum_{j=0}^{n_{\infty}-1} \xi_{j}\, z^{-j-1}+O(z^{-n_{\infty}-1}),
\end{equation}

\begin{proposition}\label{prop_unique}
If $p+n_{\infty}=2M$, then there exists at most one rational function $R$ of the form \eqref{eqn:Phi_partial_fraction} satisfying the conditions \eqref{eqn:Phi_at_z_j} and \eqref{eqn:Phi_at_infty}. Moreover, if such a rational function $R$ exists, it is necessarily real.
\end{proposition}

\begin{proof}
Suppose there exist two rational functions $R_1$ and $R_2$ of the form \eqref{eqn:Phi_partial_fraction} that satisfy \eqref{eqn:Phi_at_z_j} and \eqref{eqn:Phi_at_infty}. Each $R_i$ can be written as $R_i(z)=P_i(z)/Q_i(z)$, where $\deg(P_i)<\deg(Q_i)=M$. The fact that both $R_1$ and $R_2$ satisfy the interpolation and asymptotic conditions implies that their difference
\[
G(z):=R_1(z)-R_2(z) = \frac{P_1(z)Q_2(z)-P_2(z)Q_1(z)}{Q_1(z)Q_2(z)}
\]
has $p$ zeros at the points $\{z_j\}_{1 \le j \le p}$ and satisfies $G(z)=O(z^{-n_{\infty}-1})$ as $z\to\infty$. Consequently, the numerator $P_1Q_2-P_2Q_1$ is a polynomial with $p$ distinct zeros, yet
\[
\deg(P_1Q_2-P_2Q_1) \le \deg(Q_1Q_2)-n_{\infty}-1 = 2M-n_{\infty}-1=p-1.
\]
Thus the polynomial must vanish identically, and we conclude that $R_1(z)\equiv R_2(z)$.  

Next, suppose that $R$ is a rational function of the form \eqref{eqn:Phi_partial_fraction} satisfying \eqref{eqn:Phi_at_z_j} and \eqref{eqn:Phi_at_infty}. Define
\[
R_3(z) := \overline{R(\overline{z})},
\]
that is, $R_3$ is obtained from $R$ by conjugating all coefficients. For every $j=1,2,\dots,p$ we then have
\[
R_3(z_j) = \overline{R(\overline{z_j})}
= \overline{F(z_{p+1-j})} = F(z_j) = R(z_j).
\]
Since $\xi_j\in \r$ and the asymptotic condition \eqref{eqn:Phi_at_infty} holds, it follows that $R(z)-R_3(z)=O(z^{-n_{\infty}-1})$ as $z\to\infty$. By the same uniqueness argument as above, we conclude that $R(z)\equiv R_3(z)$ for all $z\in \c$. In particular, $R(z)\in \r$ whenever $z\in \r$, which proves that $R$ is a real rational function.
\end{proof}

\end{document}
